\section{拟射影簇的维数}\label{sec:dimension-of-quasiprojective-varieties}

本节把维数的概念从仿射簇的情形推广到拟射影簇的情形, 然后利用维数的局部性质, 把仿射情形的关键结果引入更一般的拟射影情形.

首先回顾, 我们曾把不可约仿射簇的维数定义为其函数域的超越次数 (第~\ref{sec:dimension-of-affine-varieties}~节). 此外, 在上一节, 我们引入有理函数, 把函数域的概念推广到拟射影情形. 从这一观点来看, 以下定义自然地推广了仿射情形中维数的定义.

\begin{Definition}{拟射影簇的维数}\leavevmode\label{def:dimension-of-quasiprojective-variety}
不可约拟射影簇 $X$ 的\emph{维数}是其有理函数域在 $K$ 上的超越次数:
\[
  \dim(X)=\operatorname{trdeg}_{K}(K(X)).
\]
非空可约拟射影簇的\emph{维数}是其不可约分支的维数的最大值.
\end{Definition}

在更一般的拟射影情形, 维数的一个本质性质是它的局部性. 特别地, 拟射影簇的维数可以在任一稠密开子集上确定, 下一个结果验证了这一点.

\begin{Proposition}{维数与稠密开子簇}\leavevmode\label{prop:dimension-determined-on-dense-open-subvarieties}
若 $X$ 为拟射影簇, $U\subseteq X$ 为稠密开子集, 则
\[
  \dim(X)=\dim(U).
\]
\end{Proposition}
\begin{Proof}
先假设 $X$ 不可约. 则对任一非空 (因而稠密的) 开子集 $U\subseteq X$, 有 $K(X)=K(U)$ (Proposition~\ref{prop:rational-functions-determined-on-dense-open-subvariety}), 所以 $\operatorname{trdeg}(K(X))=\operatorname{trdeg}(K(U))$, 从而 $\dim(X)=\dim(U)$.

为证明一般情形, 设 $X=X_{1}\cup\cdots\cup X_{k}$ 为 $X$ 的不可约分解. 若 $U\subseteq X$ 开且稠密, 则
\[
  U=(X_{1}\cap U)\cup\cdots\cup(X_{k}\cap U)
\]
是 $U$ 的不可约分解 (Exercise~\ref{ex:irreducible-components-of-dense-open-subvariety}). 由于每个 $X_{i}$ 不可约, 且 $X_{i}\cap U$ 为 $X_{i}$ 的非空开子集, 以上论证蕴含 $\dim(X_{i})=\dim(X_{i}\cap U)$. 因此 $X$ 的不可约分支的最大维数与 $U$ 的不可约分支的最大维数相同, 从而 $\dim(X)=\dim(U)$.
\end{Proof}

Proposition~\ref{prop:dimension-determined-on-dense-open-subvarieties} 提供了一种有用的方法, 将拟射影簇的维数理论归结到仿射情形: 若要知道某个不可约拟射影簇 $X$ 的维数, 只须选取一个仿射开子集 $U\subseteq X$ 并计算 $U$ 的维数. 下面几个例子利用了这一策略.

\begin{Example}{射影空间的维数}\leavevmode\label{ex:dimension-of-projective-space}
由于 $\mathbb{P}^{n}$ 包含 $\mathbb{A}^{n}_{0}=\mathbb{A}^{n}$ 作为稠密开子集, Proposition~\ref{prop:dimension-determined-on-dense-open-subvarieties} 蕴含
\[
  \dim(\mathbb{P}^{n})=\dim(\mathbb{A}^{n})=n.
\]
也可直接由 Definition~\ref{def:dimension-of-quasiprojective-variety} 以及 Example~\ref{ex:projective-function-field-via-affine-patch} 中 $K(\mathbb{P}^{n})=K(x_{1},\ldots,x_{n})$ 的观察得到这个结论.
\end{Example}

\begin{Example}{射影超曲面的维数}\leavevmode\label{ex:dimension-of-projective-hypersurfaces}
令 $f\in K[x_{0},\ldots,x_{n}]$ 为不可约齐次多项式, 考察不可约射影簇
\[
  X=\mathcal{V}_{\mathbb{P}}(f)\subseteq\mathbb{P}^{n}.
\]
我们想计算 $X$ 的维数. 不妨假设 $f\ne x_{0}$, 从而
\[
  X_{0}=\mathcal{V}_{\mathbb{P}}(f)\cap\mathbb{A}^{n}_{0}\ne\varnothing.
\]
因为 $X_{0}$ 是 $X$ 的稠密开子集, Proposition~\ref{prop:dimension-determined-on-dense-open-subvarieties} 蕴含 $X$ 和 $X_{0}$ 具有相同的维数. 而且, 由于 $X_{0}$ 是仿射簇, 可以利用为研究仿射情形中的维数而建立的工具. 特别地, 由于
\[
  X_{0}=\mathcal{V}_{\mathbb{A}}(f_{0})
  \qquad\text{其中}\qquad
  f_{0}(x_{1},\ldots,x_{n})=f(1,x_{1},\ldots,x_{n}),
\]
维数论基本定理 (Theorem~\ref{thm:fundamental-theorem-of-dimension-theory}) 蕴含 $X_{0}$ 的维数为 $n-1$, 因而也有 $\dim(X)=n-1$.

此外, 任一非常数齐次多项式 (不必不可约) 的消没集是其不同不可约因子的消没集的并. 因此, 可以去掉 $f$ 不可约的假设, 得到射影空间中单个非常数齐次多项式的消没集的每个不可约分支的余维都是一.
\end{Example}

\begin{Example}{Grassmann 簇的维数}\leavevmode\label{ex:dimension-of-grassmannians}
考察第~\ref{sec:grassmannians-and-plucker-maps}~节引入的 Grassmann 簇 $\mathrm{G}(k,n)\subseteq\mathbb{P}^{\binom{n}{k}-1}$. 在 Exercise~\ref{ex:grassmannian-affine-charts-and-irreducibility} 中, 要求读者证明 $\mathrm{G}(k,n)$ 不可约, 且其所有仿射限制都同构于 $\mathbb{A}^{k(n-k)}$. 因此, $\mathrm{G}(k,n)$ 包含一个同构于 $k(n-k)$ 维仿射空间的稠密开子簇, 所以
\[
  \dim(\mathrm{G}(k,n))=k(n-k).
\]
为了更直观地理解这个维数, 回顾 $\mathrm{G}(k,n)$ 的元素可以用各行线性无关的 $k\times n$ 矩阵表示. 此外, 在允许换基的意义下, 总可以选取包含一个等于单位矩阵的 $k\times k$ 子矩阵的矩阵表示. 由于剩余的 $k(n-k)$ 个矩阵项可以独立变化, 可见 $\mathrm{G}(k,n)$ 有 $k(n-k)$ 个独立参数, 与上面算出的维数一致.
\end{Example}

为了进一步说明在局部仿射情形下工作的用处, 现在建立拟射影情形中维数论基本定理的推广. 为了以最一般的形式陈述它, 引入局部超曲面的概念. 这是一种新的闭子簇, 局部可由一个正则函数的消没描述. 在下一个定义中, 对任意正则函数 $F\in K[U]$, 采用记号 $\mathcal{V}(F)=\{a\in U\mid F(a)=0\}$. 这是对迄今使用的记号 $\mathcal{V}(-)$ 的小幅而自然的推广.

\begin{Definition}{局部超曲面}\leavevmode\label{def:local-hypersurface}
令 $X$ 为拟射影簇. 若每个 $p\in X$ 都包含在某个开集 $U\subseteq X$ 中, 使 $Y\cap U=\mathcal{V}(F)$, 其中正则函数 $F\in K[U]$ 不在 $U$ 的任何不可约分支上消没, 则称非空子集 $Y\subseteq X$ 为 $X$ 中的\emph{局部超曲面}.
\end{Definition}

Definition~\ref{def:local-hypersurface} 的要点在于, 它为考察拟射影簇 $X\subseteq\mathbb{P}^{n}$ 的两类子簇提供了统一框架: 一类形如 $Y=X\cap\mathcal{V}(F)$, 其中正则函数 $F\in K[X]$ 不在 $X$ 的任何不可约分支上消没; 另一类形如 $Y=X\cap\mathcal{V}(f)$, 其中齐次多项式 $f\in K[x_{0},\ldots,x_{n}]$ 不在 $X$ 的任何不可约分支上消没.

注意, 局部超曲面是 $X$ 的闭子簇. 为了说明这一点, 注意 $\mathcal{V}(F)\subseteq U$ 是单点 $\{0\}$ 在正则映射 $F:U\to\mathbb{A}^{1}$ 下的原像, 因而由正则映射的连续性 (Theorem~\ref{thm:regular-maps-are-zariski-continuous}), 它在 $U$ 中闭. 然后由 Lemma~\ref{lem:closedness-is-local} 得局部超曲面是闭的.

现在陈述这一情形下的维数论基本定理.

\begin{Theorem}{维数论基本定理}\leavevmode\label{thm:fundamental-dimension-theorem-for-quasiprojective-varieties}
令 $X$ 为各不可约分支维数相同的拟射影簇, $Y\subseteq X$ 为 $X$ 中的局部超曲面. 则 $Y$ 的每个不可约分支的维数均为 $\dim(X)-1$.
\end{Theorem}
\begin{Proof}
令 $Z$ 为 $Y$ 的一个不可约分支, $p\in Z$. 由局部超曲面的定义, 存在包含 $p$ 的开子集 $U\subseteq X$ 和正则函数 $F\in K[U]$, 使 $F$ 不在 $U$ 的任何不可约分支上消没, 且
\[
  Y\cap U=\mathcal{V}(F).
\]
必要时限制到一个包含 $p$ 的更小开子集, 可进一步假设 $U$ 是仿射簇. 由于 $Z\cap U$ 不可约, 它必包含在 $U$ 的某个不可约分支中; 将这个分支记为 $V$. 则 $V$ 为不可约仿射簇, 且 $Z\cap V$ 是 $Y\cap V=\mathcal{V}(F|_{V})$ 的一个不可约分支. 由于 $\mathcal{V}(F|_{V})$ 包含 $p$, 且 $F|_{V}\ne0$, $\mathcal{V}(F|_{V})$ 既非空集也不是整个 $V$. 因此仿射情形中维数论基本定理的强形式蕴含
\[
  \dim(Z\cap V)=\dim(V)-1.
\]
$Z$ 的不可约性给出 $\dim(Z)=\dim(Z\cap V)$, 而 $V$ 的不可约性意味着它在 $X$ 的某个分支中稠密. $X$ 的分支的维数都相同, 所以 $\dim(V)=\dim(X)$. 从而 $\dim(Z)=\dim(X)-1$, 如所断言.
\end{Proof}

反复应用上一结果, 得到以下关于拟射影簇内消没集的维数的上界.

\begin{Corollary}{消没集的维数}\leavevmode\label{cor:dimension-bound-for-quasiprojective-vanishing-set}
令 $X\subseteq\mathbb{P}^{n}$ 为各不可约分支维数相同的拟射影簇, $Y=\mathcal{V}(f_{1},\ldots,f_{k},F_{1},\ldots,F_{\ell})\subseteq X$, 其中 $f_{1},\ldots,f_{k}\in K[x_{0},\ldots,x_{n}]$ 齐次, 且 $F_{1},\ldots,F_{\ell}\in K[X]$. 若 $Y\ne\varnothing$, 则
\[
  \dim(Y)\ge\dim(X)-k-\ell.
\]
\end{Corollary}
\begin{Proof}
Exercise~\ref{ex:dimension-bound-for-quasiprojective-vanishing-set}.
\end{Proof}

此外, 与仿射情形完全相同, 维数论基本定理给出维数的以下另一种刻画. 在许多拓扑学情形中, 它被用作维数的定义.

\begin{Corollary}{包含链与维数}\leavevmode\label{cor:dimension-is-length-of-irreducible-closed-chain}
若 $X$ 为拟射影簇, 则 $\dim(X)$ 等于使 $X$ 存在不可约闭子簇 $X_{0},\ldots,X_{d}\subseteq X$ 且满足
\[
  X_{0}\subsetneq X_{1}\subsetneq\cdots\subsetneq X_{d}\subseteq X
\]
的最大 $d$.
\end{Corollary}
\begin{Proof}
Exercise~\ref{ex:dimension-equals-maximal-length-of-chain}.
\end{Proof}

本节最后专门研究齐次多项式方程组的解. 回顾, 研究 $\mathbb{P}^{n}$ 的最初动机之一是找回多项式方程组可能\textquotedbl{}丢在无穷远处\textquotedbl{}的解. 例如, 虽然两条直线在 $\mathbb{A}^{2}$ 中可能不相交, 但我们已经看到它们在 $\mathbb{P}^{2}$ 中总相交. 实质上, $\mathbb{A}^{2}$ 中两条平行直线在 $\mathbb{P}^{2}$ 中的一个\textquotedbl{}无穷远点\textquotedbl{}相交. 下一个结果大幅推广这一观察: $\mathbb{P}^{n}$ 中 $k$ 个超曲面的交的维数总是至少为 $n-k$; 特别地, 只要 $k\le n$, 这个交就一定非空.

\begin{Proposition}{齐次方程组的维数}\leavevmode\label{prop:dimension-of-homogeneous-systems}
若 $f_{1},\ldots,f_{k}\in K[x_{0},\ldots,x_{n}]$ 为非常数齐次多项式且 $k\le n$, 则 $\mathcal{V}_{\mathbb{P}}(f_{1},\ldots,f_{k})\ne\varnothing$, 且
\[
  \dim(\mathcal{V}_{\mathbb{P}}(f_{1},\ldots,f_{k}))\ge n-k.
\]
\end{Proposition}

可以把 Proposition~\ref{prop:dimension-of-homogeneous-systems} 的陈述与仿射情形的 Corollary~\ref{cor:dimension-lower-bound-by-number-of-equations} 的陈述比较. 两个结果的关键区别是, 射影情形可以去掉仿射情形作为假设所需的非空性.

\begingroup
\renewcommand{\AJProofBegin}{%
  \par\addvspace{\topsep}\noindent
  {\sffamily\bfseries\CJKfamily{pffont}Proof of Proposition~\ref{prop:dimension-of-homogeneous-systems}}\enspace\ignorespaces
}
\begin{Proof}
若 $\mathcal{V}_{\mathbb{P}}(f_{1},\ldots,f_{k})\ne\varnothing$, 则 Corollary~\ref{cor:dimension-bound-for-quasiprojective-vanishing-set} 蕴含
\[
  \dim(\mathcal{V}_{\mathbb{P}}(f_{1},\ldots,f_{k}))\ge\dim(\mathbb{P}^{n})-k=n-k.
\]
因此只须证明 $\mathcal{V}_{\mathbb{P}}(f_{1},\ldots,f_{k})\ne\varnothing$. 为此考察仿射消没集
\[
  \mathcal{V}_{\mathbb{A}}(f_{1},\ldots,f_{k})\subseteq\mathbb{A}^{n+1}.
\]
由于 $f_{1},\ldots,f_{k}$ 是非常数齐次多项式, 它们都在原点消没. 因此 $\mathcal{V}_{\mathbb{A}}(f_{1},\ldots,f_{k})\ne\varnothing$, 并且由 Corollary~\ref{cor:dimension-lower-bound-by-number-of-equations} 得
\[
  \dim(\mathcal{V}_{\mathbb{A}}(f_{1},\ldots,f_{k}))\ge n+1-k>0.
\]
由于正维的簇必包含多于一个点, $f_{1},\ldots,f_{k}$ 必在某个非零点 $(a_{0},\ldots,a_{n})\in\mathbb{A}^{n+1}$ 消没. 这意味着 $\mathcal{V}_{\mathbb{P}}(f_{1},\ldots,f_{k})$ 包含点 $[a_{0}:\cdots:a_{n}]\in\mathbb{P}^{n}$. 因此 $\mathcal{V}_{\mathbb{P}}(f_{1},\ldots,f_{k})\ne\varnothing$, 如所需.
\end{Proof}
\endgroup

\BookExerciseHeading
\begin{BookExercises}
  \item\label{ex:quasiprojective-dimension-monotone-and-strict-for-proper-closed-subvariety} 利用维数的局部性, 将以下结果归结为仿射情形下适当的结果, 从而证明它们.
  \begin{enumerate}
    \item 若 $X,Y\subseteq\mathbb{P}^{n}$ 为拟射影簇且 $X\subseteq Y$, 则
    \[
      \dim(X)\le\dim(Y).
    \]
    \item 若进一步有 $Y$ 不可约, 且 $X\subsetneq Y$ 在 $Y$ 中闭, 则
    \[
      \dim(X)<\dim(Y).
    \]
  \end{enumerate}
  \item\label{ex:vanishing-ideal-of-pure-codimension-one-projective-variety-is-principal} 令 $X\subseteq\mathbb{P}^{n}$ 为射影簇, 其每个不可约分支的维数均为 $n-1$. 证明 $\mathcal{I}(X)$ 为主理想, 由 $\mathcal{I}(X)$ 中任一个次数最小的齐次多项式生成.
  \item\label{ex:dimension-bound-for-quasiprojective-vanishing-set} 证明 Corollary~\ref{cor:dimension-bound-for-quasiprojective-vanishing-set}.
  \item\label{ex:dimension-equals-maximal-length-of-chain} 证明 Corollary~\ref{cor:dimension-is-length-of-irreducible-closed-chain}.
  \item\label{ex:dense-image-has-no-larger-dimension-than-domain} 令 $F:X\to Y$ 为拟射影簇之间的正则映射, 使 $F(X)$ 在 $Y$ 中稠密. 证明 $\dim(Y)\le\dim(X)$. (提示: 证明拉回给出单射 $F^{*}:K(Y)\to K(X)$; 见 Exercise~\ref{ex:pullback-of-rational-functions-under-dominant-map}.)
  \item\label{ex:dimension-of-quasiprojective-product} 令 $X,Y$ 为拟射影簇. 证明
  \[
    \dim(X\times Y)=\dim(X)+\dim(Y).
  \]
  \item\label{ex:polynomial-map-from-higher-projective-space-is-constant} 证明当 $n>m$ 时, 每个多项式映射 $\mathbb{P}^{n}\to\mathbb{P}^{m}$ 都是常值的.
  \item\label{ex:regular-map-from-higher-projective-space-is-constant} 证明当 $n>m$ 时, 每个正则映射 $\mathbb{P}^{n}\to\mathbb{P}^{m}$ 都是常值的.
\end{BookExercises}
